\section{Preliminaries}\label{sec:prelim}

We fix the formalism, which follows \CBa{} and \CBb{} with two small generalisations: actions may
have negative preconditions (as in \DB{}), and the structural restrictions of the hierarchy are
the standard ones of \textcite{alford:et:al:icaps-15}.

\subsection{States and nondeterministic actions}

Let $F$ be a finite set of \emph{facts}. A \emph{state} is a subset $s\subseteq F$. An
\emph{action} is a triple $a=(\pre^+,\pre^-,\eff)$ of a positive precondition $\pre^+\subseteq F$,
a negative precondition $\pre^-\subseteq F$ and a finite, non-empty set $\eff$ of
\emph{effects} $(\add,\del)$ with $\add,\del\subseteq F$. The action is \emph{applicable} in $s$
if $\pre^+\subseteq s$ and $\pre^-\cap s=\emptyset$; its possible \emph{outcomes} are then
\[
  \Out(a,s) \;=\; \{\, (s\setminus\del)\cup\add \;\mid\; (\add,\del)\in\eff \,\}.
\]
So, when an effect adds and deletes the same fact, the addition wins. An action with a single
effect is \emph{deterministic}. \CBa{} and \CBb{} allow positive preconditions only; their
actions are the special case $\pre^-=\emptyset$. All our upper bounds hold with negative
preconditions, and so do all lower bounds. We say explicitly which reductions need them. Those that
do not are lower bounds for the original fragment as well.

The requirement $\eff\neq\emptyset$ is left implicit in the literature but matters: an action
without effects would have no outcomes, and strong solution concepts would then hold vacuously
after it.

\subsection{Task networks, methods and decomposition}

Let $A$ be a set of \emph{primitive} task names, each with an action $\delta(a)$, and $C$ a set of
\emph{compound} task names. A \emph{task network} $\tn=(T,\prec,\alpha)$ consists of a finite set
$T$ of task identifiers, a strict partial order ${\prec}$ on $T$ and a labelling
$\alpha:T\to A\cup C$. A task $t$ is \emph{unconstrained} (or \emph{first}) if no task of $T$
precedes it. Task networks are considered up to \emph{isomorphism}, i.e.\ renaming of
identifiers that preserves order and labels. A network is \emph{primitive} if all its labels are
primitive, and \emph{totally ordered} if ${\prec}$ is total.

A \emph{method} $m=(c,\tn_m)$ with $c\in C$ refines the compound name $c$ into the network
$\tn_m$. \emph{Decomposing} a task $t$ with $\alpha(t)=c$ by $m$ replaces $t$ by a fresh copy of
$\tn_m$. Every task that preceded $t$ now precedes all new tasks, and every task that $t$
preceded is now preceded by all new tasks. \emph{Progression} by an unconstrained primitive task $t$
removes $t$ from the network; it is the network-level counterpart of executing $\alpha(t)$.

A \emph{FOND HTN problem} $\Prob=(F,A,C,\delta,\M,\sI,\tnI)$ consists of the facts, the task
names with their actions, a finite set $\M$ of methods, an initial state $\sI$ and an initial
task network $\tnI$. As in the cited works there is no state goal: the objective is to refine
$\tnI$ completely and execute the resulting actions. A state goal $G$ can be expressed by
appending a task whose action has precondition $G$.

\subsection{Structural classes}

The following restrictions are standard. They constrain the methods (and possibly $\tnI$) only.

\begin{itemize}
  \item \PRIM: $\tnI$ is primitive (so methods play no role).
  \item \TO: $\tnI$ and all method networks are totally ordered. Then every network reachable by
    decomposition is totally ordered as well.
  \item \ACYC: the graph on $C$ with an edge $c\to c'$ whenever a method for $c$ contains a task
    labelled $c'$ is acyclic. Every refinement then terminates after boundedly many steps.
  \item \TR\ (tail-recursive): there is a total preorder (a \emph{stratification}) on $C$ such
    that for every method $(c,\tn_m)$ and every compound task $t$ of $\tn_m$, $\alpha(t)$ lies in
    a strictly lower stratum than $c$, unless $t$ is the \emph{last} task of $\tn_m$ (all other
    tasks of $\tn_m$ precede it), in which case $\alpha(t)$ lies in the same or a lower stratum.
    Nothing is required of $\tnI$.
  \item \REG\ (regular): every method network contains at most one compound task, and if it does,
    that task is last. \textcite{bercher:et:al:ijcai-19} impose this on methods only;
    \textcite{erol:et:al:amai-96} and \CBa{} also require it of $\tnI$. We write \REGtn{} for the
    latter class. Regular problems are tail-recursive with a single stratum.
  \item \ARB: no restriction.
\end{itemize}

\CBa{} and \CBb{} phrase the acyclic and tail-recursive classes with stratifications whose
definitions are garbled in places (a ``total order'' with ties, an orientation of acyclicity that
contradicts reflexivity, and an ambiguous notion of ``last''). The definitions above are those of
\textcite{alford:et:al:icaps-15} and are what all published proofs actually use. The distinction
between \REG{} and \REGtn{} looks innocent but matters a great deal for strong outcome-dependent
solutions (Section~\ref{sec:fixed-reg}).

\subsection{Complexity and computability}

Problems are given by explicit descriptions: actions with their lists of effects, methods,
initial state and initial network, written down in binary. Plan existence for a solution
concept and a class is the set of descriptions in the class that have a solution. We use the
standard classes $\PTIME\subseteq\NP\subseteq\PSPACE\subseteq\EXPTIME\subseteq\NEXPTIME\subseteq
\EXPSPACE\subseteq\TWOEXP$, polynomial-time many-one reductions, and the alternation theorems of
\textcite{chandra:et:al:jacm-81}: alternating polynomial time is \PSPACE, alternating polynomial
space is \EXPTIME, alternating exponential time is \EXPSPACE, and alternating exponential space is
\TWOEXP. By Savitch's theorem \parencite{savitch:jcss-70}, nondeterministic and deterministic
space coincide.

Most of our upper bounds are proved by exhibiting a \emph{game}. An AND/OR game is played on
configurations; at existential configurations the planner picks a successor, at universal
configurations the environment does. If configurations have polynomial (exponential) size and
plays have polynomial (exponential) length, deciding the game is in \PSPACE{} (\EXPSPACE{}); if
configurations have polynomial (exponential) size but plays may be long, an attractor computation
over all configurations gives \EXPTIME{} (\TWOEXP). Fair games, which we need for strong cyclic
solutions, have a nested fixpoint (Section~\ref{sec:sc}), and the same bounds apply to them.

For undecidable problems we use the following notions. A set of words is \emph{recursively
enumerable} (\RE, or semi-decidable) if some program halts exactly on its elements, and
\emph{co-recursively enumerable} (\coRE) if its complement is \RE. A set $L$ is in
$\Sigma^1_1$ if there is a computable relation $R$ such that
\[
  w\in L \iff \exists f:\mathbb N\to\mathbb N\ \forall n:\ R\bigl(w, f(0)\cdots f(n-1)\bigr),
\]
i.e.\ membership is witnessed by an infinite object that can be checked on its finite prefixes
(Kleene's normal form, see \textcite{rogers:67}). Every arithmetical set is reducible to a
$\Sigma^1_1$-complete set but not conversely, so $\Sigma^1_1$-complete sets are not arithmetical.

\paragraph{Two-stack machines.}
Several undecidability proofs use \emph{two-stack machines} \parencite{minsky:67}: a finite
control with two stacks over a finite alphabet. Each transition is applicable in one control
state, pops a given symbol from one of the two stacks, pushes a word onto the same stack and moves
to a new control state. Two-stack machines simulate Turing machines step by step (the two stacks
hold the tape left and right of the head), so their halting problem is undecidable. Moreover,
deterministic Turing machines are simulated by deterministic two-stack machines, which we use in
Section~\ref{sec:unrestricted}.
